% Cover letter: Brightwell Analytics, Senior Data Engineer (BW-DE-0417).
\documentclass[11pt,a4paper]{article}
\usepackage{lmodern}
\usepackage[T1]{fontenc}
\usepackage[utf8]{inputenc}
\usepackage[margin=2.5cm]{geometry}
\usepackage[hidelinks]{hyperref}
\pagestyle{empty}
\setlength{\parindent}{0pt}
\setlength{\parskip}{8pt}

\begin{document}

\textbf{Alex Fernsby}\\
\href{mailto:alex.fernsby@example.com}{alex.fernsby@example.com} \quad +1 555 0100

\vspace{6pt}
Hiring Manager, Brightwell Analytics\\
Reference BW-DE-0417

\vspace{6pt}
Dear Hiring Manager,

I am applying for the Senior Data Engineer role in your data platform team. I have 7 years of experience
building batch and streaming pipelines for retail and logistics companies, and two of the strongest parts
of the job you describe are already close to what I do. You want batch pipelines over retail data; I spent
nearly four years building them at Marlow Retail Group, including the company's first Airflow deployment.
You want nightly pipelines that feed forecasting and analytics and are fast enough to fix when they break;
at Tidewater Logistics I moved the nightly delivery-performance mart to incremental dbt models and cut its
rebuild from 6 hours to 55 minutes.

On the core requirements, Python and SQL are my strongest tools. At Tidewater I own a pipeline moving 40
million Kafka events per day into PostgreSQL and BigQuery under Airflow, and Great Expectations tests on 30
core tables cut failed loads by 70 percent over two quarters. The retail side adds a customer-orders star
schema used by 12 analysts, Python ingestion services for partner price feeds at about 2 million rows a day,
and support for a warehouse move to BigQuery. Your nice-to-haves fit as well: Kafka is a core part of my
current pipeline, and I mentor two junior engineers and lead our weekly data-modelling review.

I want to be straightforward about four requirements I do not meet fully. Scala and Spark are not yet
learned; they are new to me, and I would expect a real learning curve on both in the first months. Most of
the pipeline logic I write runs on Python and SQL today, and I do not claim any of it as Scala or Spark
experience. Second, my testing is real but not universal: Great Expectations
runs on 30 core tables, not every table. Third, my data-quality work is test-based; I cannot point to
production monitoring built the way you describe it. Fourth, my deployments so far have not involved
working with data scientists to productionise features. Each of these is work I would want to take on
deliberately in the first months rather than assume I already do.

I would welcome the chance to talk about how your team works. I am based in Rotterdam and open to hybrid
roles in the Netherlands, which fits your two days a week on site.

Kind regards,\\
Alex Fernsby

\end{document}
